% Options for packages loaded elsewhere
\PassOptionsToPackage{unicode}{hyperref}
\PassOptionsToPackage{hyphens}{url}
\documentclass[
  a4paper,
]{article}
\usepackage{xcolor}
\usepackage[margin=25mm]{geometry}
\usepackage{amsmath,amssymb}
\setcounter{secnumdepth}{-\maxdimen} % remove section numbering
\usepackage{iftex}
\ifPDFTeX
  \usepackage[T1]{fontenc}
  \usepackage[utf8]{inputenc}
  \usepackage{textcomp} % provide euro and other symbols
\else % if luatex or xetex
  \usepackage{unicode-math} % this also loads fontspec
  \defaultfontfeatures{Scale=MatchLowercase}
  \defaultfontfeatures[\rmfamily]{Ligatures=TeX,Scale=1}
\fi
\usepackage{lmodern}
\ifPDFTeX\else
  % xetex/luatex font selection
\fi
% Use upquote if available, for straight quotes in verbatim environments
\IfFileExists{upquote.sty}{\usepackage{upquote}}{}
\IfFileExists{microtype.sty}{% use microtype if available
  \usepackage[]{microtype}
  \UseMicrotypeSet[protrusion]{basicmath} % disable protrusion for tt fonts
}{}
\makeatletter
\@ifundefined{KOMAClassName}{% if non-KOMA class
  \IfFileExists{parskip.sty}{%
    \usepackage{parskip}
  }{% else
    \setlength{\parindent}{0pt}
    \setlength{\parskip}{6pt plus 2pt minus 1pt}}
}{% if KOMA class
  \KOMAoptions{parskip=half}}
\makeatother
\usepackage{longtable,booktabs,array}
\newcounter{none} % for unnumbered tables
\usepackage{calc} % for calculating minipage widths
% Correct order of tables after \paragraph or \subparagraph
\usepackage{etoolbox}
\makeatletter
\patchcmd\longtable{\par}{\if@noskipsec\mbox{}\fi\par}{}{}
\makeatother
% Allow footnotes in longtable head/foot
\IfFileExists{footnotehyper.sty}{\usepackage{footnotehyper}}{\usepackage{footnote}}
\makesavenoteenv{longtable}
\setlength{\emergencystretch}{3em} % prevent overfull lines
\providecommand{\tightlist}{%
  \setlength{\itemsep}{0pt}\setlength{\parskip}{0pt}}
\usepackage{newunicodechar}
\newunicodechar{↔}{\ensuremath{\leftrightarrow}}
\newunicodechar{→}{\ensuremath{\rightarrow}}
\newunicodechar{←}{\ensuremath{\leftarrow}}
\newunicodechar{≥}{\ensuremath{\ge}}
\newunicodechar{≤}{\ensuremath{\le}}
\newunicodechar{×}{\ensuremath{\times}}
\newunicodechar{≪}{\ensuremath{\ll}}
\newunicodechar{≫}{\ensuremath{\gg}}
\newunicodechar{∈}{\ensuremath{\in}}
\newunicodechar{∞}{\ensuremath{\infty}}
\newunicodechar{±}{\ensuremath{\pm}}
\newunicodechar{−}{\ensuremath{-}}
\newunicodechar{≈}{\ensuremath{\approx}}
\newunicodechar{≠}{\ensuremath{\ne}}
\newunicodechar{∝}{\ensuremath{\propto}}
\newunicodechar{√}{\ensuremath{\sqrt{\phantom{x}}}}
\newunicodechar{⊕}{\ensuremath{\oplus}}
\newunicodechar{α}{\ensuremath{\alpha}}
\newunicodechar{β}{\ensuremath{\beta}}
\newunicodechar{θ}{\ensuremath{\theta}}
\newunicodechar{ρ}{\ensuremath{\rho}}
\newunicodechar{π}{\ensuremath{\pi}}
\newunicodechar{φ}{\ensuremath{\varphi}}
\newunicodechar{λ}{\ensuremath{\lambda}}
\newunicodechar{σ}{\ensuremath{\sigma}}
\newunicodechar{Δ}{\ensuremath{\Delta}}
\newunicodechar{Σ}{\ensuremath{\Sigma}}
\newunicodechar{ν}{\ensuremath{\nu}}
\newunicodechar{κ}{\ensuremath{\kappa}}
\newunicodechar{μ}{\ensuremath{\mu}}
\newunicodechar{τ}{\ensuremath{\tau}}
\newunicodechar{ω}{\ensuremath{\omega}}
\newunicodechar{ε}{\ensuremath{\varepsilon}}
\newunicodechar{Ψ}{\ensuremath{\Psi}}
\newunicodechar{⁻}{\ensuremath{{}^{-}}}
\newunicodechar{¹}{\ensuremath{{}^{1}}}
\newunicodechar{²}{\ensuremath{{}^{2}}}
\newunicodechar{³}{\ensuremath{{}^{3}}}
\newunicodechar{⋆}{\ensuremath{\star}}
\newunicodechar{∗}{\ensuremath{*}}
\newunicodechar{★}{\ensuremath{\bigstar}}
\newunicodechar{∎}{\ensuremath{\blacksquare}}
\newunicodechar{ℝ}{\ensuremath{\mathbb{R}}}
\newunicodechar{ℤ}{\ensuremath{\mathbb{Z}}}
\newunicodechar{₀}{\ensuremath{{}_{0}}}
\newunicodechar{₁}{\ensuremath{{}_{1}}}
\newunicodechar{₂}{\ensuremath{{}_{2}}}
\newunicodechar{₃}{\ensuremath{{}_{3}}}
\newunicodechar{₄}{\ensuremath{{}_{4}}}
\newunicodechar{₅}{\ensuremath{{}_{5}}}
\newunicodechar{₆}{\ensuremath{{}_{6}}}
\newunicodechar{₇}{\ensuremath{{}_{7}}}
\newunicodechar{₈}{\ensuremath{{}_{8}}}
\newunicodechar{₉}{\ensuremath{{}_{9}}}
\usepackage{bookmark}
\IfFileExists{xurl.sty}{\usepackage{xurl}}{} % add URL line breaks if available
\urlstyle{same}
\hypersetup{
  hidelinks,
  pdfcreator={LaTeX via pandoc}}

\author{}
\date{}

\begin{document}

\section{Lock Dynamics in an Anonymous Two-Channel Closed Wave
System}\label{lock-dynamics-in-an-anonymous-two-channel-closed-wave-system}

\subsection{Demonstration of the Reality of the Quantization Mechanism
(Arnold Tongue) and an Elimination Program for Order-Selection
Mechanisms --- Measured Tongue Widths, the Fixed-Point Map, and Clock
Inheritance of Observational
Selection}\label{demonstration-of-the-reality-of-the-quantization-mechanism-arnold-tongue-and-an-elimination-program-for-order-selection-mechanisms-measured-tongue-widths-the-fixed-point-map-and-clock-inheritance-of-observational-selection}

\textbf{Author:} Noriaki Kihara\\
\textbf{ORCID:} 0009-0004-6753-4020\\
\textbf{Date:} August 3, 2026\\
\textbf{Version DOI:} \texttt{10.5281/zenodo.21764000}\\
\textbf{Concept DOI:} \texttt{10.5281/zenodo.21763999}\\
\textbf{Position:} Additional paper v1 of the two-channel
exchange-scattering system in the ``Wave Information Readout'' series
(ninth-paper series; \textbf{Part III of a trilogy}. Part I =
Two-Grammar Decomposition {[}A{]}; Part II = Structure of Counting
Readouts {[}B{]})

\begin{center}\rule{0.5\linewidth}{0.5pt}\end{center}

\subsection{Abstract}\label{abstract}

\textbf{(Zeroth-order no-go)} The standard readout of the anonymous
two-channel closed wave system (phase-blind per-bin power ratios) is an
exact constant of the motion (drift \(\le2.2\times10^{-16}\); theorem in
{[}A{]}), and neither dynamics of the coupling angle \(\theta\) nor
locking exists.

\textbf{(Reality of the quantization mechanism)} Under a phase-sensitive
readout (the power ratio of the symmetric channel \(X=A+B\), introduced
as a conservative working hypothesis), \(\theta\) evolves and \textbf{a
genuine Arnold tongue is real}: over a finite interval of initial
amplitudes \([2.518,\ \ge3.44]\) (width \(\ge0.92\)), the mean of the
late-time winding number locks exactly to \(\theta/\pi=1/3\) (machine
precision over 800 collisions; \(\theta\) itself oscillates on a
periodic orbit --- rational mean winding). Under the root convention,
\((m,n)=(1,6)\), order 6. Together with the locks of orders 3, 4, and 6
obtained by the counting construction ({[}B{]}, the complete menu of the
crystallographic restriction), \textbf{the mechanism of quantization
plateaus is real in an endogenously coupled system in which the coupling
is not externally imposed}.

\textbf{(Elimination of selection mechanisms)} Locking, however, is not
a mechanism of \textbf{selection}. (i) In-model measurement of tongue
widths: against the width \(\ge0.92\) for denominator 3, tongues with
denominator \(\ge4\) in the search band have width \(<0.002\) ---
\textbf{a width ratio \(>461\)}. The weight of denominator 124 is
negligible. (ii) Immediately below the lower edge of the tongue lies a
supercritical region in which the winding number varies
non-monotonically; the closest approach to the target \(\rho=39/124\) is
\(3.3\times10^{-4}\), with no plateau detected. (iii) The attractors of
the second dynamical candidate, the fixed-point-type readout (cross
spectrum), are distributed over a continuous band \([0.016,0.450]\) with
\textbf{no rational concentration} (0/41) --- eliminated. (iv)
Observational selection (repeated observation confirming recurrence to
the initial state) yields fidelity exactly 1 at exact finite-order
roots, directly showing that \textbf{selection is inherited by the
divisor structure of the observation clock}: an observer with clock
\(J=8\) cannot see the order-124 recurrence (\(F_8=0.0026\)), while an
observer with \(J=248\) can (\(F_{248}=1.000000000000001\)). However,
the fidelity landscape is smooth in the neighborhood of the roots, and
finitely many observations can narrow the state down only to the
recurrence \textbf{basin} (concentration ratio \(1.001\) at \(n=20\),
recorded as a refutation of prediction P2).

\textbf{(Consequence)} Quantization (why discrete) is explained by
locking; selection (why this value) is not explained by locking. With
the elimination program complete, the order-124 selection problem has
been transformed into a single falsifiable question: \textbf{``why is
the observation clock commensurable with 248?''} This final form is
taken up by the parent paper (the main paper of the ninth-paper series).

\textbf{Keywords:} Arnold tongue, mode locking, devil's staircase,
supercritical circle map, return fidelity, observational selection,
quantum recurrence, reproducible computation

\begin{center}\rule{0.5\linewidth}{0.5pt}\end{center}

\subsection{0. Conclusion}\label{conclusion}

\[
\boxed{
\begin{aligned}
&\text{Under the phase-blind readout, no }\theta\text{ dynamics exists (no-go, }2.2\times10^{-16}\text{).}\\
&\text{Under the phase-sensitive readout, an Arnold tongue is real: }\rho=1/3\text{ exact, width}\ge0.92\text{ (mechanism of quantization).}\\
&\text{By the width ratio}>461\text{, fixed-point elimination, and clock inheritance of observational selection,}\\
&\text{neither locking nor fixed points select order 124.}\\
&\text{Final form of the selection problem: "why is the observation clock commensurable with 248?"}
\end{aligned}
}
\]

\subsection{Position of This Paper (Three-Part
Structure)}\label{position-of-this-paper-three-part-structure}

Part I {[}A{]} establishes the binomial identity decomposition of the
flow and the two grammars (addition/product, sign, neutralization,
hierarchy); Part II {[}B{]} establishes the structure of counting
readouts (meta accounting rules, convention audit, Niven intersections,
integer-pair readouts). This paper (Part III) treats dynamics: the
reality of the mechanism of quantization plateaus, and an elimination
program for candidate order-selection mechanisms. The three parts share
a single reproduction package.

\subsection{1. Research Problem}\label{research-problem}

By {[}B{]}, the locks reachable by the counting construction are limited
to the Niven intersections (orders 3, 4, 6). Meanwhile, the root
\(R_{124,23}\) corresponding to the physical \(\alpha\) has order 124
{[}2{]}, and this selection lies outside counting. The central problem
of the ninth-paper series is ``what selects order 124?'', and the
candidates are organized into four: (a) static state structure (already
refuted in {[}B{]}), (b) dynamical locking, (c) fixed-point-type
attractors, (d) selection by observation (recurrence confirmation). This
paper adjudicates (b), (c), and (d) by experiment.

\subsection{2. System and Design
Boundary}\label{system-and-design-boundary}

The system, the state construction, the standard readout, and the
rotation are identical to Section 2 of {[}A{]} (scattering core and
readout unchanged; target values are used only for state construction).
The introduction of the phase-sensitive readout is made explicit as a
conservative working hypothesis (Section 4), and the ``observation'' of
observational selection is defined as the evaluation of the return
fidelity to the initial state pair,

\[
F_J=\frac{|\langle a_0|a_J\rangle+\langle b_0|b_J\rangle|^2}{(\langle a_0|a_0\rangle+\langle b_0|b_0\rangle)^2}
\]

(no back-action on the dynamics; a pure readout).

\subsection{3. Zeroth-Order No-Go}\label{zeroth-order-no-go}

The standard readout is a function only of the per-bin rotation
invariants \(|A_k|^2+|B_k|^2\) and is exactly conserved under real
orthogonal rotations (theorem in {[}A{]}, Section 4). Numerical
confirmation: over 400 collisions × 4 amplitudes, the maximum drift of
\(\theta\) is \(2.2\times10^{-16}\). Hence dynamics, locking, and
selection of \(\theta\) do not exist at zeroth order. All of the
dynamics below rests on an explicit choice of a readout that breaks this
no-go.

\subsection{4. Readout Classification (Dynamical
Aspect)}\label{readout-classification-dynamical-aspect}

The seven \(\theta\) readouts were classified by their behavior under
the dynamics they themselves drive (the dynamical aspect of the table in
{[}A{]}, Section 4). Only the diagonal sum is conserved. Within the
dynamical class, the cross-spectrum readout (R4) is of
\textbf{fixed-point-convergent type} (it moves substantially away from
the initial value and then freezes; initially misidentified as conserved
and corrected by the classification experiment --- see the Addendum),
while the \(X\)-power readout (V1) is of \textbf{oscillatory, locking
type}. In what follows, V1 is used for the lock probe and R4 for the
fixed-point probe.

\subsection{5. Reality of the Arnold Tongue (Quantization
Mechanism)}\label{reality-of-the-arnold-tongue-quantization-mechanism}

\subsubsection{5.1 Discovery and
Refinement}\label{discovery-and-refinement}

Sweeping the initial B amplitude under V1, there exists a finite
interval in which the tail average of the late-time winding number
sticks exactly to a rational number:

\begin{itemize}
\tightlist
\item
  A coarse scan (61 points × 400 collisions) detected a plateau, and a
  fine scan (13 points × 800 collisions, tail 300) confirmed it:
  \textbf{for amplitudes 2.6--3.3, \(\overline{\theta/\pi}=1/3\) to
  machine precision (distance 0.0, maximum \(5.6\times10^{-17}\))}
\item
  The lower edge was fixed at 2.518 by a 39/124 micro-zoom (27 points ×
  1500 collisions), and the upper edge at \(\ge3.44\) by an additional
  scan --- \textbf{width \(\ge0.92\)}
\item
  \(\theta\) itself oscillates (tail standard deviation
  \(\sim7\times10^{-2}\)), and only the mean winding number locks
  rationally on a periodic orbit --- the standard behavior of mode
  locking in circle maps
\end{itemize}

Under the root convention \(\theta=\pi/2-\pi m/n\), \((m,n)=(1,6)\).
Together with the counting locks of {[}B{]} (orders 3, 4, 6), the
mechanism of quantization plateaus in an endogenously coupled system is
established.

\subsubsection{5.2 In-Model Measurement of Tongue Widths (First Negation
of
Selection)}\label{in-model-measurement-of-tongue-widths-first-negation-of-selection}

Not a single tongue with denominator \(\ge4\) was detected in the search
band \([2.470, 2.522]\) (resolution 0.002, 1500 collisions). Therefore,
as an in-model measurement,

\[
\frac{\mathrm{width}(q=3)}{\mathrm{width}(q\ge4)}>\frac{0.92}{0.002}=461
\]

and the width collapses rapidly with the denominator. \textbf{The tongue
weight of denominator 124 is negligible on our own data, without even
invoking external theory (the general theory of circle maps).}

\subsubsection{5.3 Supercritical Diagnosis and the 39/124 Upper Bound
(Second
Negation)}\label{supercritical-diagnosis-and-the-39124-upper-bound-second-negation}

Immediately below the lower edge of the 1/3 tongue, the winding number
varies non-monotonically by \(\sim10^{-2}\) between adjacent amplitude
points --- the symptom of a supercritical region in which tongues
overlap. The closest approach to the target
\(\rho=39/124=0.3145161\ldots\) is \(3.3\times10^{-4}\) (amplitude
2.504, 1500 collisions), with no plateau detected. A naive sweep search
for high-order tongues is in principle unsuitable with this readout, and
the coupling-interpolation campaign originally designed became
\textbf{unnecessary for the conclusion} owing to the measured ratio of
5.2 (the history is recorded in the Addendum).

\subsection{6. Fixed-Point Map (Elimination of the Third
Candidate)}\label{fixed-point-map-elimination-of-the-third-candidate}

The attractor \(\theta^*\) under the self-driven R4 readout was measured
at 41 initial-amplitude points (400 collisions). The attractors are
distributed over a continuous band \(\theta^*/\pi\in[0.016, 0.450]\),
with \textbf{no exact concentration on rational angles whatsoever}
(0/41). Fixed-point-type attractors select no special value --- the
third candidate is eliminated.

\subsection{7. Observational Selection (Fourth Candidate: The Mechanism
Is Real, Selection Is Inherited by the
Clock)}\label{observational-selection-fourth-candidate-the-mechanism-is-real-selection-is-inherited-by-the-clock}

\subsubsection{7.1 Design and Predictions}\label{design-and-predictions}

For a state family (61 amplitude-family points plus 2 exact-root states:
\(R=1/2\) and \(R_{124,23}\), constructed by inverse search), we measure
the return fidelity \(F_J\) at clocks \(J\in\{8,248\}\) and the survival
weight \(F_J^n\) after \(n=20\) repeated observations. Predictions
(fixed before measurement): P1 --- \(F=1\) only at exact roots
(\(R=1/2\) at both clocks by \(8\mid248\); the 124 root only at
\(J=248\)); P2 --- concentration ratio onto recurrent states \(>10^3\)
at \(n=20\).

\subsubsection{7.2 Results}\label{results}

\textbf{P1 holds in full}: \(F_8(R{=}1/2)=0.999999999999999\),
\(F_{248}(\text{124 root})=1.000000000000001\), and \textbf{clock
selectivity} --- \(F_8(\text{124 root})=0.0026\). That is, \textbf{which
recurrences are ``visible'' is determined by the divisor structure of
the observation clock}: an observer with \(J=8\) is blind in principle
to the order-124 recurrence, while an observer with \(J=248\) selects
both orders 8 and 124. Selection is inherited not by the state but by
the observation clock.

\textbf{P2 is refuted} (recorded as designed): the concentration ratio
at \(n=20\) is \(1.001\). The fidelity landscape is smooth in the
neighborhood of the roots, and neighboring states also survive with
\(F\approx1-\varepsilon\), so that \textbf{finitely many observations
can select the recurrence basin, but the exact point only
asymptotically}. This is a quantitative limit of observational
selection, consistent with the statement in the published paper {[}2{]}
that ``the residual vanishes in ideal arithmetic, so the depth is
formally infinite.''

\subsubsection{7.3 Consequence: Coordinate Transformation of the
Selection
Problem}\label{consequence-coordinate-transformation-of-the-selection-problem}

The adjudication of the four candidates is complete: static structure
(no --- {[}B{]}), dynamical locking (no --- width ratio 461), fixed
points (no --- continuous band), observational selection (the mechanism
is real, but selection is inherited by the commensurability \(P\mid J\)
of the observation clock). Hence ``why order 124?'' has been transformed
into the single question

\[
\boxed{\text{"Why is the observation clock commensurable with 248?"}}
\]

This is not a retreat but an advance: the question has moved from a
search of the state space to a problem of the structure of the observer
(register, resolution). This final form is taken up by the parent paper
(the main paper of the ninth-paper series: the generation structure of
fermions and wave-packet collapse).

\subsection{8. Claims}\label{claims}

\textbf{Claim 1.} Under the phase-blind readout, no dynamics of the
coupling angle exists (no-go; theorem plus numerics
\(2.2\times10^{-16}\)).

\textbf{Claim 2.} Under the phase-sensitive readout, an Arnold tongue is
real: \(\rho=1/3\) exact, width \(\ge0.92\), \((m,n)=(1,6)\). Together
with the counting locks (orders 3, 4, 6 {[}B{]}), \textbf{the mechanism
of quantization (why discrete) is real in an endogenously coupled
system}.

\textbf{Claim 3.} Locking is not a mechanism of selection (why this
value): by in-model measurement, tongue widths collapse by a factor
\(>461\) between denominator 3 and denominator ≥4, and the weight of
order 124 is negligible.

\textbf{Claim 4.} Fixed-point-type channels do not select (continuous
band of attractors; rational concentration 0/41).

\textbf{Claim 5.} Observational selection is real as a mechanism
(fidelity 1 at exact roots) and \textbf{inherits selection into the
divisor structure of the observation clock} (direct demonstration of
clock selectivity). However, finitely many observations can narrow the
state down only to the recurrence basin (P2 refuted). Hence the final
form of the selection problem is ``why is the observation clock
commensurable with 248?''

\subsection{9. Relation to Known
Systems}\label{relation-to-known-systems}

Mode locking, the devil's staircase, and the denominator dependence of
tongue widths are established mathematics of circle maps {[}13{]}, and
quantization plateaus produced by phase locking exist physically as
superconducting Shapiro steps {[}12{]}. The contributions of this paper
are: (i) demonstrating locking in an endogenously coupled system in
which the coupling is read from the state rather than being an apparatus
constant; (ii) completing the elimination program for selection
candidates by measuring tongue widths, fixed points, and observational
selection in the same system; (iii) directly showing the inheritance of
selection into the divisor structure of the observation clock.

\subsection{10. Open Problems}\label{open-problems}

\begin{enumerate}
\def\labelenumi{\arabic{enumi}.}
\tightlist
\item
  \textbf{Commensurability of the observation clock}: the final form of
  ``why 248?''. Model the observer itself as a closure (register) within
  the system, and search for a mechanism by which its intrinsic period
  becomes commensurable with the dominant recurrence period --- passed
  to the parent paper
\item
  \textbf{Derivation of the V1 readout}: the phase-sensitive readout is
  a working hypothesis. The conserved-readout classification ({[}A{]},
  Section 4) narrowed the candidate space, but there is no first
  principle for the readout that the dynamics selects
\item
  \textbf{Staircase in the subcritical region}: measurement of the
  complete staircase via the coupling interpolation \(P_\lambda\) became
  unnecessary for the conclusion (5.2), but remains as a quantitative
  theory of tongue-width scaling
\end{enumerate}

All of the above are open problems internal to the model. The overall
picture of the problems that remain as connections to observational
physics is organized collectively, as a separate series, in the next
section.

\subsection{11. Physical-Connection Problems Beyond the Trilogy --- From
the Three Directions to 3+1-Dimensional Spacetime, and from the Two
Grammars to Gravity and the Coulomb
Force}\label{physical-connection-problems-beyond-the-trilogy-from-the-three-directions-to-31-dimensional-spacetime-and-from-the-two-grammars-to-gravity-and-the-coulomb-force}

Since this is the final paper of the trilogy, we conclude by making
explicit, in a single problem table, the range established inside the
model and the range that remains as connection to observational physics.

First, we distinguish the logical hierarchy. The \textbf{internal
selection problem} left by the preceding sections (``why is the
observation clock commensurable with 248?'') asks why a particular
finite-order root is selected inside the model. The \textbf{physical
connection problems} of this section ask how the internal structure thus
selected is realized as the observed 3+1-dimensional spacetime, gravity,
and the Coulomb force. The two are non-competing, separate series of
problems.

There is no need to rediscover what has already been derived: the
spontaneous emergence of the three directions (the metastabilization
process reported by the earlier stage of this series and by the parent
paper), the inverse-square law via harmonic closure {[}3{]}, the unique
differentiation of the gravity-type and charge-type two grammars
({[}A{]} Section 6.5, exclusion of the inverted assignment), and the
structure of counting, finite order, locking, and the observation clock
({[}B{]} and this paper). What remains is the construction of the
\textbf{realization map} that carries these internal structures onto
observational physics --- concretely, the following problems.

\textbf{(1) Local spatialization of the three directions.} Let
\(T_1,T_2,T_3\) be the generating actions of the three directions that
arise at metastability, and show that they form the basis of a local
three-dimensional space rather than remaining three internal modes. The
zeroth-order candidate is the Gram metric
\(g_{ij}(\Psi)=\mathrm{Re}\langle T_i\Psi|T_j\Psi\rangle\); verify
whether \(\mathrm{rank}\,g=3\) and positive definiteness hold for the
family of metastable states. No physical meaning may be assigned to the
numbering of the three directions (anonymity); the requirement is that
the space be definable as equivalent under basis rotations \(SO(3)\).
\textbf{The existence of three directions and the establishment of a
three-dimensional space are not the same thing} --- this is the core of
what must be examined throughout this section.

\textbf{(2) Globalization of the local three directions and
integrability.} Do the three directions obtained at each closure and at
each local state form a common space? Examine the structure of
\([T_i,T_j]\), the parallel transport of local bases and its path
dependence (holonomy), and Frobenius-type integrability conditions, and
determine whether the local three directions can be connected into a
single global space.

\textbf{(3) The time direction and the construction of 3+1-dimensional
spacetime.} This system has multiple time candidates: collision count,
phase advance, finite-order recurrence, and the observation clock \(J\).
Determine which of these (or which combination) carries physical time,
and verify whether, together with the three positive-definite
directions, a Lorentz-type metric
\(ds^2=-c_{\mathrm{eff}}^2 d\tau^2+g_{ij}dx^i dx^j\), causal structure,
and a finite propagation speed arise endogenously.

\textbf{(4) Identity of distances.} Derive the identity of the
harmonic-closure distance, the geometric distance constructed from the
three directions, and the distance read by the observer,

\[
r_{\mathrm{harmonic}}=r_{\mathrm{geometric}}=r_{\mathrm{observed}}
\]

(the completed form of the \(\Delta\theta\leftrightarrow r\) dictionary
of {[}A{]} Open Problem 2). A scalar \(r\) alone does not suffice: one
must construct the direction \(\hat{\mathbf r}\) joining two closures
and a gradient operation, and show that a \textbf{directed}
inverse-square field corresponding to
\(\nabla(1/r)=-\hat{\mathbf r}/r^2\) can be defined on the emergent
three directions.

\textbf{(5) Map from the gravity-type grammar to physical gravity.} Map
the phase-blind, additive source readout of the magnitude term onto
physical mass-energy sources, and formalize how sources deform the
emergent metric and connection (magnitude-type source →
\(\delta g_{\mu\nu}\) or \(\delta\Gamma^\mu{}_{\nu\rho}\) → motion of
closures). The first verification targets are: Newton-type acceleration
in the weak-field, low-velocity limit; universal free fall independent
of the internal structure of the test closure; agreement of the
inertial-side and gravitational-side readouts; and the endogenous
derivation of the hierarchy exponent linking the coupling to the scale
ratio \(R/R_0\) (the scale dictionary of {[}A{]}).

\textbf{(6) Map from the charge-type grammar to the physical Coulomb
force.} Read the overlap term
\(\mathrm{Re}\langle a|b\rangle\sim|q_A||q_B|\cos(\phi_B-\phi_A)\) as a
charge-type source vertex, connect it to the mediation by harmonic
closure {[}3{]} and to the response on the receiving side, and derive
\(\mathbf F_{AB}=K\,q_Aq_B\,\hat{\mathbf r}_{AB}/r_{AB}^2\) on the
emergent space (the structure source vertex × mediator × response
vertex). Items to be settled individually: whether the sign of the
relative phase maps to approach or recession (Bjerknes type or
electromagnetic type --- the automatic-adjudication prediction of
{[}A{]}); whether the neutralization of non-shared channels maps to
physical electrical neutrality; the conservation law of charge-type
sources and a Gauss-type flux law; the path by which
\(1-R_\ast=\sin^2(23\pi/124)\) maps, unit conventions included, to the
observed elementary charge; and whether the introduction of
time-dependent phases produces magnetic and radiative components beyond
the static Coulomb field.

\textbf{(7) A spacetime dictionary common to the two interactions.}
Separate distances and separate spaces must not be assigned to the
gravity type and the charge type. On the same emergent metric, the same
\(r\), and the same \(\tau\), the magnitude-type grammar → gravitational
response and the overlap-type grammar → charge response must hold
\textbf{simultaneously}, and this correspondence must be invariant or
covariant under basis rotations of the three directions, A/B channel
exchange, renaming of harmonics, changes of readout convention, and the
observer's choice of coordinates.

\textbf{(8) Falsification conditions for the physical connection map.}
The connection is not to be certified on grounds of similarity. If any
of the following holds, the identification with physical gravity and the
Coulomb force is rejected: no stable rank-3 metric can be constructed
from the three directions / the local three directions are not globally
integrable / the harmonic-closure distance and the emergent geometric
distance do not agree / magnitude-type sources produce no universal
geometric response / the phase sign of the overlap type does not map to
consistent attraction and repulsion / the same charge magnitude does not
appear on both the source side and the response side / the two grammars
do not act simultaneously on the same 3+1-dimensional spacetime.

The above are not problems that return the internal results of this
trilogy to an unsettled state; they are the next-stage connection
problems of mapping the derived pre-geometric structures onto
observables. The current position of this series is organized as
follows:

\[
\boxed{
\begin{aligned}
&\text{Spontaneous emergence of the three directions} && \text{derived (earlier stage of this series)}\\
&\text{Inverse-square mediation} && \text{derived [3]}\\
&\text{Gravity-type and charge-type two grammars (unique assignment)} && \text{derived and verified [A]}\\
&\text{Structure of counting, rational quantization, and the observation clock} && \text{derived, limits verified ([B], this paper)}\\
&\text{Three directions}\to 3+1\text{-dimensional spacetime} && \text{unconnected (Problems 1–4)}\\
&\text{Two grammars}\to\text{physical gravity and Coulomb force} && \text{unconnected (Problems 4–7)}
\end{aligned}
}
\]

The conclusion of the trilogy closes in three steps: \textbf{The two
grammars differentiate spontaneously, and the assignment is unique
({[}A{]}). Their readout carries the structure of counting, finite
order, locking, and the observation clock ({[}B{]}, this paper). The
realization map to physical spacetime and physical forces is taken up,
as the eight problems above, by the next stage --- the parent paper and
its successors.}

\subsection{12. Reproducibility}\label{reproducibility}

The runners for all experiments, all measurement CSVs, and the figures
are committed, following the same discipline as {[}A{]} (SHA-verified
core references, pre-measurement fixing of predictions, anchor
reproduction, and preservation of refutations and corrections).

{\def\LTcaptype{none} % do not increment counter
\begin{longtable}[]{@{}
  >{\raggedright\arraybackslash}p{(\linewidth - 4\tabcolsep) * \real{0.3333}}
  >{\raggedright\arraybackslash}p{(\linewidth - 4\tabcolsep) * \real{0.3333}}
  >{\raggedright\arraybackslash}p{(\linewidth - 4\tabcolsep) * \real{0.3333}}@{}}
\toprule\noalign{}
\begin{minipage}[b]{\linewidth}\raggedright
Experiment
\end{minipage} & \begin{minipage}[b]{\linewidth}\raggedright
Folder
\end{minipage} & \begin{minipage}[b]{\linewidth}\raggedright
Commit
\end{minipage} \\
\midrule\noalign{}
\endhead
\bottomrule\noalign{}
\endlastfoot
Mode-locking probe (Z₆ discovery, no-go) &
\texttt{mode\_locking\_probe\_pre\_v1} & 56453653 \\
1/3-tongue fine scan & same as above,
\texttt{run\_one\_third\_tongue\_fine\_scan\_v1} & 15ee20cb \\
Readout classification + staircase + 39/124 zoom &
\texttt{tongue\_spectroscopy\_pre\_v1} & 913756fe, 15ee20cb \\
Tongue-width measurement (upper edge, width ratio 461) &
\texttt{tongue\_width\_measurement\_v1} & 47e2da02 \\
Fixed-point map & \texttt{partial\_sharing\_fixed\_point\_v1} (E7 part)
& 47e2da02 \\
Observational selection (clock selectivity) &
\texttt{observation\_selection\_v1} & 47e2da02 \\
\end{longtable}
}

\begin{center}\rule{0.5\linewidth}{0.5pt}\end{center}

\section{References}\label{references}

\subsection{Self-citations}\label{self-citations}

\begin{enumerate}
\def\labelenumi{\arabic{enumi}.}
\tightlist
\item
  Noriaki Kihara, ``Anonymous Equal-Amplitude Composite-Wave Model:
  Basic Axiom System v6,'' Version DOI:
  \texttt{10.5281/zenodo.21465984}, Concept DOI:
  \texttt{10.5281/zenodo.21315735}, 2026.
\item
  Noriaki Kihara, ``Discovery of Finite-Order Resonance in Iterated
  Exchange Scattering --- Identifying the Origin of Peaks near
  Fine-Structure-Constant Values 137 and 128 with a Reproducible
  Wave-Packet Model,'' Version DOI: \texttt{10.5281/zenodo.21421367},
  Concept DOI: \texttt{10.5281/zenodo.21421366}, 2026.
\item
  Noriaki Kihara, ``Future Phase-Position Acceleration Map and the
  Inverse-Square Law via Harmonic Closure in an AB Two-Body Closed Phase
  System v4,'' Version DOI: \texttt{10.5281/zenodo.21468270}, Concept
  DOI: \texttt{10.5281/zenodo.21441081}, 2026. (Precedent for the
  three-step procedure: zeroth-order no-go → conservative working
  hypothesis → independent verification)
\item
  Noriaki Kihara, ``Two-Grammar Decomposition of Interaction in an
  Anonymous Two-Channel Closed Wave System'' (Part I of the trilogy),
  Version DOI: \texttt{10.5281/zenodo.21763996}, Concept DOI:
  \texttt{10.5281/zenodo.21763995}, 2026.
\item
  Noriaki Kihara, ``Structure of Counting Readouts in an Anonymous
  Two-Channel Closed Wave System'' (Part II of the trilogy), Version
  DOI: \texttt{10.5281/zenodo.21763998}, Concept DOI:
  \texttt{10.5281/zenodo.21763997}, 2026.
\item
  Noriaki Kihara, ``Preliminary Summary of Acceleration-Basis and
  Localization Exchange in Exchange-Interference Scattering-Matrix
  Fermion-Like Collisions v1,'' Version DOI:
  \texttt{10.5281/zenodo.21333768}, 2026. (Code provenance of the
  engine)
\item
  Noriaki Kihara, ``Selection of the State-Exchange Weight G\_R = 1 − R
  and Candidate Correspondence with the Fine-Structure Constant: A
  Numerical Experiment v1,'' Version DOI:
  \texttt{10.5281/zenodo.21396761}, 2026. (Code provenance of System A)
\end{enumerate}

\subsection{External References}\label{external-references}

\begin{enumerate}
\def\labelenumi{\arabic{enumi}.}
\setcounter{enumi}{11}
\tightlist
\item
  S. Shapiro, \emph{Phys. Rev.~Lett.} \textbf{11}, 80 (1963). DOI:
  \texttt{10.1103/PhysRevLett.11.80}.
\item
  M. H. Jensen, P. Bak, and T. Bohr, \emph{Phys. Rev.~Lett.}
  \textbf{50}, 1637 (1983). DOI: \texttt{10.1103/PhysRevLett.50.1637}.
\end{enumerate}

\begin{center}\rule{0.5\linewidth}{0.5pt}\end{center}

\textbf{Addendum (record of corrections and history)} (i) The
cross-spectrum readout R4 was misidentified as ``a second conserved
readout'' → corrected after the classification experiment showed it to
be of fixed-point-convergent type. (ii) The concentration prediction P2
of observational selection (\(>10^3\) at \(n=20\)) was refuted ---
recorded in the main text as a principled limit of finitely many
observations due to the smooth fidelity landscape. (iii) The originally
designed subcritical-staircase campaign via coupling interpolation
(\(P_\lambda\)) became unnecessary for the conclusion because the
in-model measurement of tongue widths (width ratio \(>461\)) gave the
same conclusion more directly. The decision history is recorded together
with the judgment.

\end{document}
